\documentclass[12pt]{article}
% Préambule commun à tous les documents extraits de « La Revue CAPES »
\usepackage[T1]{fontenc}
\usepackage[utf8]{inputenc}
\usepackage{lmodern}
\usepackage{amsmath,amssymb,amsthm,amsfonts,mathrsfs}
\usepackage{setspace}
\usepackage{listings}
\usepackage{pgfplots}
\pgfplotsset{compat=1.18}
\usepackage{longtable}
\usepackage{adjustbox}
\usepackage{lettrine}
\usepackage{tkz-tab}
\usepackage{changepage}
\usepackage{pdflscape}
\usepackage{graphicx}
\usepackage{tikz}
\usepackage{xcolor}
\usepackage[a4paper,top=2cm,bottom=2.5cm,left=2.5cm,right=2cm]{geometry}
\usepackage{fancyhdr}
\usepackage{draftwatermark}
\SetWatermarkText{La RevueCapes}
\SetWatermarkLightness{0.9}
\SetWatermarkScale{2}
\usepackage{hyperref}
\hypersetup{colorlinks=true,linkcolor=blue!50!black,urlcolor=blue!50!black}

\definecolor{revue}{RGB}{20,60,120}

\pagestyle{fancy}
\fancyhf{}
\renewcommand{\headrulewidth}{0.4pt}
\setlength{\headheight}{15pt}

% \entete{Matière}{Titre}{Type de document}
\newcommand{\entete}[3]{%
  \fancyhead[L]{\small\textsc{La Revue CAPES} -- #1}%
  \fancyhead[R]{\small #2}%
  \fancyfoot[C]{\thepage}%
  \thispagestyle{plain}%
  \begin{center}
    {\color{revue}\rule{\textwidth}{1.2pt}}\\[4pt]
    {\small\textsc{Concours directs CAP-CEG / CAPES -- Burkina Faso}}\\[6pt]
    {\LARGE\bfseries\color{revue} #2}\\[6pt]
    {\large\itshape #1 -- #3}\\[2pt]
    {\color{revue}\rule{\textwidth}{1.2pt}}
  \end{center}
  \vspace{0.5cm}%
}

% Encadré pour les remarques ajoutées dans les corrigés
\newcommand{\remarque}[1]{\par\medskip\noindent\fcolorbox{revue}{revue!6}{\parbox{\dimexpr\linewidth-2\fboxsep-2\fboxrule}{\textbf{\color{revue}Remarque.} #1}}\par\medskip}


\begin{document}
\entete{Mathématiques}{Corrigé détaillé du sujet type n°1}{Correction}
\begin{spacing}{1.5}

\medskip\textbf{Exercice 1}

\textbf{1. Montrons que $a$ divise $b-a$ et déterminons le quotient.}

Comme $a$ divise $b$, il existe un entier $k$ tel que $b=ka$. En retranchant $a$ aux deux membres de cette égalité, on obtient :
\[
b-a=ka-a=(k-1)a .
\]
Comme $k-1$ est un entier, $a$ divise $b-a$, et le quotient de $b-a$ par $a$ est $k-1=\dfrac{b}{a}-1$ (lorsque $a\neq 0$).

De plus, comme $0\le a\le b$, on a $k\ge 1$ (si $a\neq0$), donc $k-1\in\mathbb{N}$ : le quotient est bien un entier naturel.

\textbf{2. Entiers naturels $n$ tels que $n-1$ divise $n+3$.}

\emph{Condition nécessaire.} Pour que $n-1$ soit un entier naturel non nul diviseur de $n+3$, il faut $n\ge 2$ (si $n=1$, $n-1=0$ ne divise pas $4$ ; si $n=0$, $n-1=-1$ n'est pas un entier naturel).
Supposons donc $n\ge 2$ et $n-1$ divise $n+3$. On a $n-1\le n+3$, donc d'après la question 1 (avec $a=n-1$ et $b=n+3$), $n-1$ divise
\[
(n+3)-(n-1)=4 .
\]
Les diviseurs positifs de $4$ sont $1$, $2$ et $4$. Donc $n-1\in\{1;2;4\}$, c'est-à-dire $n\in\{2;3;5\}$.

\emph{Réciproque (vérification).} Si $n-1$ divise $4$, alors $n-1$ divise $(n-1)+4=n+3$. On vérifie :
$n=2$ : $1$ divise $5$ ; \quad $n=3$ : $2$ divise $6$ ; \quad $n=5$ : $4$ divise $8$.

\textbf{Conclusion :} les entiers naturels $n$ tels que $n-1$ divise $n+3$ sont $\boxed{2,\ 3 \text{ et } 5}$.

\medskip\textbf{Exercice 2}

\textbf{1. Matrice $M(a)$ de $f_a$ dans la base canonique.}

Les colonnes de $M(a)$ sont les images des vecteurs de la base canonique : $f_a(e_1)=(2-a,\,2(1-a),\,0)$, $f_a(e_2)=(a-1,\,2a-1,\,0)$ et $f_a(e_3)=(0,0,a)$. D'où
\[
M(a)=\begin{pmatrix}
2-a & a-1 & 0 \\
2(1-a) & 2a-1 & 0\\
0 & 0 & a
\end{pmatrix}.
\]

\textbf{2. Noyaux $K(0)$ et $K(1)$.}

Par définition, $K(a)=\{(x,y,z)\in\mathbb{R}^3 \ / \ f_a(x,y,z)=(0,0,0)\}$, c'est-à-dire l'ensemble des solutions du système
\[
(S_a)\ :\ \begin{cases}
(2-a)x+(a-1)y=0\\
2(1-a)x+(2a-1)y=0\\
az=0
\end{cases}
\]

\underline{Cas $a=0$.} Le système devient $2x-y=0$, $2x-y=0$, $0=0$. Donc $y=2x$ et $z$ est quelconque :
\[
K(0)=\{(x,2x,z)\ /\ x,z\in\mathbb{R}\}=\operatorname{Vect}\big((1,2,0),(0,0,1)\big).
\]
C'est un plan vectoriel (de dimension 2), de base $\{(1,2,0);(0,0,1)\}$.

\underline{Cas $a=1$.} Le système devient $x=0$, $y=0$, $z=0$ (en fait $f_1=\mathrm{Id}_{\mathbb{R}^3}$ car $M(1)=I_3$). Donc
\[
K(1)=\{(0,0,0)\}.
\]

\textbf{3. Noyau $K(a)$ pour $a\neq0$ et $a\neq1$.}

La troisième équation $az=0$ donne $z=0$ puisque $a\neq 0$.

Les deux premières équations forment un système linéaire homogène en $(x,y)$ dont le déterminant vaut
\[
\Delta=(2-a)(2a-1)-2(1-a)(a-1)=(-2a^2+5a-2)+2(a-1)^2=-2a^2+5a-2+2a^2-4a+2=a .
\]
Comme $a\neq 0$, $\Delta\neq 0$ : ce système (de Cramer) n'admet que la solution nulle $x=y=0$. Par conséquent
\[
K(a)=\{(0,0,0)\}\quad\text{pour tout } a\neq 0 .
\]
\remarque{On retrouve le cas $a=1$. Plus généralement, $\det M(a)=a\cdot\Delta=a^2$, donc $f_a$ est bijective si et seulement si $a\neq0$ : son noyau est alors réduit au vecteur nul.}

\textbf{4. Écriture $M(a)=A+aB$.}

On sépare dans chaque coefficient la partie constante et la partie proportionnelle à $a$ :
\[
M(a)=\begin{pmatrix}
2 & -1 & 0 \\
2 & -1 & 0 \\
0 & 0 & 0
\end{pmatrix}
+a\begin{pmatrix}
-1 & 1 & 0 \\
-2 & 2 & 0 \\
0 & 0 & 1
\end{pmatrix},
\qquad\text{donc}\qquad
A=\begin{pmatrix}
2 & -1 & 0 \\
2 & -1 & 0 \\
0 & 0 & 0
\end{pmatrix},\quad
B=\begin{pmatrix}
-1 & 1 & 0 \\
-2 & 2 & 0 \\
0 & 0 & 1
\end{pmatrix}.
\]
(On remarque que $A=M(0)$.)

\textbf{5. Calcul de $AB$, $BA$, $A^2$, $B^2$ et de $M(a)^2$.}

\[
AB=\begin{pmatrix}
2(-1)+(-1)(-2) & 2(1)+(-1)(2) & 0\\
2(-1)+(-1)(-2) & 2(1)+(-1)(2) & 0\\
0&0&0
\end{pmatrix}=\begin{pmatrix}0&0&0\\0&0&0\\0&0&0\end{pmatrix},
\]
\[
BA=\begin{pmatrix}
-2+2 & 1-1 & 0\\
-4+4 & 2-2 & 0\\
0&0&0
\end{pmatrix}=\begin{pmatrix}0&0&0\\0&0&0\\0&0&0\end{pmatrix},
\]
\[
A^2=\begin{pmatrix}
4-2 & -2+1 & 0\\
4-2 & -2+1 & 0\\
0&0&0
\end{pmatrix}=A,
\qquad
B^2=\begin{pmatrix}
1-2 & -1+2 & 0\\
2-4 & -2+4 & 0\\
0&0&1
\end{pmatrix}=B .
\]
Ainsi $AB=BA=0$, $A^2=A$ et $B^2=B$ ($A$ et $B$ sont des matrices de projecteurs).

En développant (le produit matriciel est distributif, mais on garde l'ordre des facteurs) :
\[
M(a)^2=(A+aB)(A+aB)=A^2+aAB+aBA+a^2B^2=A+a^2B .
\]

\textbf{6. Calcul de $M(a)^n$.}

Montrons par récurrence que, pour tout $n\in\mathbb{N}$, $\mathcal{P}(n)$ : $M(a)^n=A+a^nB$.

\emph{Initialisation.} Pour $n=0$ : $A+B=\begin{pmatrix}1&0&0\\0&1&0\\0&0&1\end{pmatrix}=I_3=M(a)^0$. $\mathcal{P}(0)$ est vraie (et $\mathcal{P}(1)$ est la question 4).

\emph{Hérédité.} Supposons $M(a)^n=A+a^nB$ pour un entier $n$. Alors
\[
M(a)^{n+1}=M(a)^n\,M(a)=(A+a^nB)(A+aB)=A^2+aAB+a^nBA+a^{n+1}B^2=A+a^{n+1}B .
\]
\emph{Conclusion.} Pour tout entier naturel $n$ :
\[
M(a)^n=A+a^nB=\begin{pmatrix}
2-a^n & a^n-1 & 0 \\
2(1-a^n) & 2a^n-1 & 0\\
0 & 0 & a^n
\end{pmatrix}.
\]
On remarque que $M(a)^n=M(a^n)$.

\medskip\textbf{Exercice 3}

\textbf{1. Calcul de $B=\displaystyle\sum_{p=0}^{n}C_n^p$.}

D'après la formule du binôme de Newton, pour tous réels $x,y$ : $(x+y)^n=\displaystyle\sum_{p=0}^{n}C_n^p x^{p}y^{n-p}$. Avec $x=y=1$ :
\[
2^n=(1+1)^n=\sum_{p=0}^{n}C_n^p\,1^{p}\,1^{n-p}=\sum_{p=0}^{n}C_n^p,\qquad\text{donc}\qquad \boxed{B=2^n}.
\]
\emph{Interprétation :} $B$ est le nombre total de parties d'un ensemble à $n$ éléments.

\textbf{2. Relation $C_n^p=f(n,p)\,C_{n-1}^{p-1}$ (pour $1\le p\le n$).}
\begin{align*}
C_n^p=\frac{n!}{p!\,(n-p)!}&=\frac{n\,(n-1)!}{p\,(p-1)!\,\big((n-1)-(p-1)\big)!}\\
&=\frac{n}{p}\cdot\frac{(n-1)!}{(p-1)!\,\big((n-1)-(p-1)\big)!}=\frac{n}{p}\,C_{n-1}^{p-1}.
\end{align*}
Donc $\boxed{f(n,p)=\dfrac{n}{p}}$, autrement dit $p\,C_n^p=n\,C_{n-1}^{p-1}$.

\textbf{3. Calcul de $S=\displaystyle\sum_{p=0}^{n}p\,C_n^p$ (pour $n\ge 1$).}

Le terme d'indice $p=0$ est nul, donc $S=\displaystyle\sum_{p=1}^{n}p\,C_n^p$. D'après la question 2, $p\,C_n^p=n\,C_{n-1}^{p-1}$, d'où
\[
S=\sum_{p=1}^{n}n\,C_{n-1}^{p-1}=n\sum_{p=1}^{n}C_{n-1}^{p-1}\underset{r=p-1}{=}n\sum_{r=0}^{n-1}C_{n-1}^{r}=n\,2^{n-1}
\]
en utilisant la question 1 avec $n-1$ à la place de $n$. Ainsi $\boxed{S=n\,2^{n-1}}$ (et $S=0$ si $n=0$).

\remarque{On peut aussi dériver $x\mapsto(1+x)^n=\sum_{p=0}^n C_n^p x^p$ : $n(1+x)^{n-1}=\sum_{p=1}^n pC_n^px^{p-1}$, puis prendre $x=1$.}

\medskip\textbf{Exercice 4}

$(E)\ :\ y''+y'+y=0$.

\textbf{1. Résolution de $(E)$ et comportement des solutions.}

$(E)$ est une équation différentielle linéaire du second ordre, homogène, à coefficients constants. Son équation caractéristique est
\[
r^2+r+1=0,\qquad \Delta=1-4=-3=(i\sqrt3)^2 .
\]
Elle admet deux racines complexes conjuguées $r_{1,2}=-\dfrac12\pm i\dfrac{\sqrt3}{2}$, c'est-à-dire $\alpha\pm i\beta$ avec $\alpha=-\dfrac12$ et $\beta=\dfrac{\sqrt3}{2}$.

Les solutions de $(E)$ sont donc les fonctions définies sur $\mathbb{R}$ par
\[
y(x)=e^{-\frac{x}{2}}\left(\lambda\cos\Big(\frac{\sqrt3}{2}x\Big)+\mu\sin\Big(\frac{\sqrt3}{2}x\Big)\right),\qquad \lambda,\mu\in\mathbb{R}.
\]
En posant $R=\sqrt{\lambda^2+\mu^2}$ et en choisissant $\varphi$ tel que $\lambda=R\cos\varphi$, $\mu=R\sin\varphi$, on peut aussi écrire
\[
y(x)=R\,e^{-\frac{x}{2}}\cos\Big(\frac{\sqrt3}{2}x-\varphi\Big).
\]

\underline{Comportement en $+\infty$.} Pour tout $x$, $|y(x)|\le R\,e^{-x/2}$ et $\lim\limits_{x\to+\infty}e^{-x/2}=0$. Par encadrement, \textbf{toute solution tend vers $0$ en $+\infty$} (oscillations amorties).

\underline{Comportement en $-\infty$.} Si $R=0$, $y$ est la fonction nulle. Si $R\neq 0$, on a $\lim\limits_{x\to-\infty}e^{-x/2}=+\infty$ : la solution oscille avec une amplitude $R\,e^{-x/2}$ qui tend vers $+\infty$. Plus précisément, pour $x_k=\dfrac{2}{\sqrt3}(\varphi+2k\pi)$, on a $y(x_k)=R\,e^{-x_k/2}$, et $x_k\to-\infty$ quand $k\to-\infty$ ; donc $y$ n'est pas bornée au voisinage de $-\infty$ et n'a pas de limite en $-\infty$ (elle prend aussi les valeurs $-R\,e^{-x/2}$ aux points $x_k+\frac{2\pi}{\sqrt3}$).

\textbf{En résumé :} la seule solution de $(E)$ bornée sur $\mathbb{R}$ est la fonction nulle.

\textbf{2. a) $g:x\mapsto f(x+T)-f(x)$ est solution de $(E)$.}

$f$ est de classe $C^2$ sur $\mathbb{R}$, donc $g$ aussi, avec $g'(x)=f'(x+T)-f'(x)$ et $g''(x)=f''(x+T)-f''(x)$. Posons $h=f''+f'+f$ ; par hypothèse $h(x+T)=h(x)$ pour tout réel $x$. Alors
\begin{align*}
g''(x)+g'(x)+g(x)&=\big[f''(x+T)+f'(x+T)+f(x+T)\big]-\big[f''(x)+f'(x)+f(x)\big]\\
&=h(x+T)-h(x)=0 .
\end{align*}
Donc $g$ est solution de $(E)$.

\textbf{b) Si $f$ est bornée sur $\mathbb{R}$, alors $f$ est $T$-périodique.}

Supposons qu'il existe $M>0$ tel que $|f(x)|\le M$ pour tout $x\in\mathbb{R}$. Alors, pour tout $x$,
\[
|g(x)|\le|f(x+T)|+|f(x)|\le 2M ,
\]
donc $g$ est une solution de $(E)$ bornée sur $\mathbb{R}$. D'après la question 1, la seule solution bornée de $(E)$ est la fonction nulle (une solution non nulle n'est pas bornée au voisinage de $-\infty$). Donc $g=0$, c'est-à-dire
\[
f(x+T)=f(x)\quad\text{pour tout } x\in\mathbb{R}:
\]
$f$ est $T$-périodique.

\remarque{Il est essentiel ici d'utiliser le comportement en $-\infty$ : en $+\infty$ toutes les solutions tendent vers $0$, ce qui ne permettrait pas de conclure.}

\end{spacing}
\end{document}
